\section*{第 \ref{ch:b3:covariant-electromagnetism} 章——协变电动力学}

\begin{solution}{exo:b3:covariant-electromagnetism:1}
(a) $\gamma = 1.25$：$a'^0 = 1.25(5 - 1.8) = 4.0$，$a'^1 = 1.25(3 -
3) = 0$；$a\cdot a = 25 - 9 = 16 = 16 - 0$。(b) 该变换是
线性的，故对求和可分配。(c) $(c, \vect v)$：否——它的
“时间”分量对所有粒子都相同，混合关系却要求不然；
$\gamma(c, \vect v) = \dd x^\mu/\dd\tau$：是，即
四维矢量除以不变量 $\dd\tau$。(d) $\vect E$ 的三个分量
与 $\vect B$ 的分量相混合，而不与任何标量混合：这些场
构成一个二指标张量，而非四维矢量。
\end{solution}

\begin{solution}{exo:b3:covariant-electromagnetism:2}
(a) $u = I/nSe = 10/(\num{8.5e28} \times \num{e-6} \times
\num{1.6e-19}) = \qty{7.4e-4}{m/s}$。(b) 晶格：$(nec, \vect 0)$；
电子：$(-nec, -ne\vect u)$ ——它们的电流密度
$\rho_-\vect u = -ne\vect u$ 与其漂移方向相反，这
\emph{恰是}常规电流的方向。(c) 两者均静止且均匀：各项皆为零。(d)
总计：$(0, \vect\jmath)$，其中 $j = I/S = \qty{e7}{A/m^2}$ ——
纯电流而无电荷，这一构型使导线的相对论效应
变得微妙。
\end{solution}

\begin{solution}{exo:b3:covariant-electromagnetism:3}
(a) $E_y$：$F^{20} = E_y/c$；$B_x$：$F^{32} = B_x$。(b) 对于 $\vect B
= B\vect e_z$ 时，仅 $F^{12} = -B$ 与 $F^{21} = +B$ 非零；对于波，
$F^{20} = E/c$，$F^{12} = -E/c$（及其反对称伙伴）。(c)
$F^{12} = \partial^1A^2 - \partial^2A^1 = -\partial_xA_y +
\partial_yA_x = -(\operatorname{\vect{curl}}\vect A)_z = -B_z$：
与矩阵一致。(d) 粒子能量变化率
$\propto F^{0\nu}u_\nu$；在静止参考系中 $u_\nu = (c, \vect 0)$，且
反对称性使 $F^{00} = 0$：这是一个永不对
静止粒子做功的力——磁场部分的定义性质。
\end{solution}

\begin{solution}{exo:b3:covariant-electromagnetism:4}
(a) 沿场方向做洛伦兹变换：纵向分量不变，$\vect
E' = E_0\vect e_y$，$\vect B' = \vect 0$。(b) 沿 $\vect
e_x$ 变换：$E'_y = \gamma E_0$，$B'_z = -\gamma vE_0/c^2$ ——极板
现在呈现为面电流 $\pm\sigma'v$，两个相反的
电流片恰好围出这样的场。(c) $E'^2 - c^2B'^2 =
\gamma^2E_0^2(1 - \beta^2) = E_0^2$；$\vect E'\cdot\vect B' = 0$：
两者均保持不变。(d) 极板沿 $x$ 方向收缩，因此 $\sigma' =
\gamma\sigma$，与 $E' = \sigma'/\varepsilon_0 = \gamma
E_0$ 一致；垂直于变换方向的间距保持不变。
\end{solution}

\begin{solution}{exo:b3:covariant-electromagnetism:5}
(a) 磁力 $-e\vect v\wedge\vect B$ 推动电子
沿杆运动：电动势 $= vBL$。(b) 在杆的静止参考系中，杆处于
静止状态——静止电荷不受磁力——但
变换给出 $\vect E' = \gamma\,\vect v\wedge\vect B$：
一个真实的电场起驱动作用。(c) $E'L = \gamma vBL
\approx vBL$，精确到 $v/c$ 的一阶。(d) 之所以必须如此，是因为“动生”与
“感生”电动势是同一现象在两个参考系中的不同描述：
磁通量法则不过是披着 1831 年外衣的协变性——爱因斯坦
1905 年的论文正是从这一磁体与导体的不对称性开始的。
\end{solution}

\begin{solution}{exo:b3:covariant-electromagnetism:6}
(a) $x$ 分量保持不变；对于横向分量，
$E_y'^2 + E_z'^2 - c^2(B_y'^2 + B_z'^2) = \gamma^2[(E_y - vB_z)^2 +
(E_z + vB_y)^2 - c^2(B_y + vE_z/c^2)^2 - c^2(B_z - vE_y/c^2)^2]$;
交叉项相消，平方项合并为 $\gamma^2(1 -
\beta^2) = 1$ 倍的未变换组合。(b) 对
$E_xB_x + E_yB_y + E_zB_z$ 做同样的整理。(c) $E^2 - c^2B^2 =
3c^2B^2 > 0$：沿 $\vect E\wedge\vect B$ 方向以 $v = c^2B/E =
c/2$ 变换；此时 $B' = 0$，留存的场为 $E' = E/\gamma =
\sqrt3\,cB$ ——其平方恰为不变量，这是必然的。(d) 光波的两个不变量
均为零：任一参考系都必须保持 $E' = cB' \perp$，而不可能是纯场。
\end{solution}

\begin{solution}{exo:b3:covariant-electromagnetism:7}
(a) 令 $v_{\text{d}} = E/B$：$E'_y = \gamma(E - v_{\text{d}}B) =
0$；$B'_z = \gamma(B - v_{\text{d}}E/c^2) = B/\gamma$：纯磁场，略有
减弱。(b) 在该系中为角频率 $qB'/\gamma m$ 的圆周运动；回到实验室系，
该圆周运动叠加匀速漂移——一条
沿与两场均垂直的方向以 $E/B$ 缓慢漂移的摆线，正是
磁控管中电荷的轨迹。(c) 以恰好 $v_{\text{d}}$ 运动的粒子
在漂移参考系中静止，此时唯一的场为磁场，它
不受任何作用，不发生偏转。(d) 若 $E > cB$，所需漂移
速度超过 $c$，因此不存在这样的参考系；取而代之的是纯 $E$ 参考系
（上一练习），电荷
沿场方向被无限加速。
\end{solution}

\begin{solution}{exo:b3:covariant-electromagnetism:8}
(a) 相位 $\omega t - \vect k\cdot\vect r$ 计数的是波峰穿越事件，
这是一个不变数；它等于 $k^\mu x_\mu$，并且
对所有 $x^\mu$ 该乘积的不变性要求 $k^\mu$
按四维矢量变换。(b) $\omega'/c = \gamma(\omega/c -
\beta k_x)$，其中 $k_x = (\omega/c)\cos\theta$；当 $\theta = 0$ 时，
$\omega' = \omega\sqrt{(1-\beta)/(1+\beta)}$。(c) $\cos\theta' =
k'_x/(\omega'/c)$：即所述公式。(d) $\beta = \num{e-4}$：
$20.5''$；一年之中，视位置扫出一个以该值为
角半轴的椭圆——布拉德利光行差。
\end{solution}

\begin{solution}{exo:b3:covariant-electromagnetism:9}
(a) 在静止参考系中为库仑场；变换使横向分量乘以
$\gamma$（在横向平面上，$E'_\perp = \gamma
q/4\pi\varepsilon_0b^2$），而沿运动方向的场，
在前方或后方、实验室距离为 $r$ 处取值，映射到静止系中的距离
为 $\gamma r$：场强减小为 $1/\gamma^2$。(b) 角宽度为
$1/\gamma$ 的饼状场以 $v$ 掠过：持续时间 $\sim (b/\gamma)/v$。(c)
$E' = 7250 \times \num{1.6e-19} \times \num{9e9}/\num{e-4} \approx
\qty{0.1}{V/m}$，持续时间 $2b/\gamma c \approx \qty{9}{fs}$。(d)
$B' = vE'/c^2 \approx E'/c$，两个不变量均为 $O(1/\gamma^2)$：
局部上与光脉冲无法区分——这是
快速电荷碰撞的“等效光子”描述的基础。
\end{solution}

\begin{solution}{exo:b3:covariant-electromagnetism:10}
(a) $E'_y = \gamma(E - vB) = \gamma E_0(1 - \beta)\cos(\cdots) =
E_0\sqrt{(1-\beta)/(1+\beta)}\cos(\cdots)$，且 $B'_z = E'_y/c$，由
同样的代数推得。(b) 相位以相同的多普勒因子变换：
仍是同一列光波，更偏红、更弱。(c) 能量密度
随多普勒因子的平方减小；光子数不变——
每个光子的 $h\nu$ 承载整个因子。(d) “
伴行”意味着因子 $\to 0$（当 $\beta \to 1$ 时）：波
该波永不变为静止，因其不变量为零——
不存在使光静止的参考系，这正是
本书相对论的起点。
\end{solution}

\begin{solution}{exo:b3:covariant-electromagnetism:11}
(a) 每圈 $\Delta E = P \times 2\pi r/c$，即题中所给结果。(b)
$\gamma = \num{e11}/\num{5.11e5} = \num{1.96e5}$: $\Delta E \approx
\qty{2.9}{GeV}$ 每圈——即束流能量的百分之三，
每圈重新注入；LEP 的速调管曾是
地球上最大的无线电发射机。(c) 质子：$\gamma = 7250$，$\gamma^4$
小 $(m_{\text{p}}/m_{\text{e}})^4 \approx \num{e13}$ 倍，约
\qty{5}{keV} 每圈——可忽略。(d) $\gamma^4 = (E/mc^2)^4$：
等能量时电子辐射多 $\num{e13}$ 倍，因此
采用直线对撞机（只辐射一次，而非每圈），或专供电子使用的巨型环形加速器，
而质子环可放心地扩大至 \qty{27}{km}。
\end{solution}

\begin{solution}{exo:b3:covariant-electromagnetism:12}
(a) $|\vect p|$ 为常数（不做功）；$\dot{\vect p} = q\vect
v\wedge\vect B$ 使 $\vect p$ 以 $qvB/p = qB/\gamma m$ 的速率转动。(b)
每圈滑移 $2\pi(\gamma - 1)$；随 $\gamma - 1$ 线性增长
至 $\Gamma$，总滑移 $\approx \pi N\Gamma$；四分之一周期为
$\pi/2$：$\Gamma \approx 1/2N$。取 $N = 100$：$\Gamma =
\num{5e-3}$，$E_k \approx \qty{4.7}{MeV}$——量级正确；实际
机器借助高 D 形盒电压，将其提升至 $\sim\qty{20}{MeV}$
（圈数减少）。(c) 同步回旋加速器：在每个脉冲期间将射频
向下扫描，以跟随 $qB/\gamma m$。等时回旋加速器：令
$B(r)$ 按 $\propto\gamma(r)$ 增长，使比值始终不变，
束流保持连续。(d) $\gamma = 1 + 590/938 = 1.63$：边缘
磁场须为中心磁场的该倍数。
\end{solution}

\begin{solution}{pb:b3:covariant-electromagnetism:1}
\textbf{1.} $u = I/nSe = \qty{7.4e-4}{m/s}$; $u^2/c^2 =
\num{6e-24}$.
\textbf{2.} $\lambda = I/u = nSe = \qty{1.36e4}{C/m}$——每束流中
每米十四千库仑。
\textbf{3.} $B = \mu_0I/2\pi d = \qty{2.0e-4}{T}$; $F = qvB =
\num{e-6} \times \num{e5} \times \num{2e-4} = \qty{2e-5}{N}$,
力的方向\emph{指向}导线（平行电流相互吸引）。
\textbf{4.} 两条线电荷恰好抵消：$\lambda_{\text{总}}
= 0$，零阶时无场。
\textbf{5.} 每米多出一个电子：$E = 2k\lambda_1/d =
\qty{2.9e-7}{V/m}$，力 $\qty{2.9e-13}{N}$——比磁力小 $10^8$ 倍。
导线每米可隐藏一亿
个杂散电子，静电力才会与磁力相抗衡：
将该力归于 $\vect B$ 是安全的。
\textbf{6.} $\qty{2e-5}{N}$ 相当于二十个沙粒的重量：
易于测量。
\textbf{7.} 晶格：$-v$。电子（以速率 $u$ 漂移，与 $I$ 方向相反，
即与变换方向相反）：速率 $(u + v)/(1 + uv/c^2)$——快于
晶格。
\textbf{8.} $\lambda' = \gamma_v(\lambda_{\text{总}} - vI/c^2) =
-\gamma_v vI/c^2$.
\textbf{9.} $\lambda' = -\num{e5} \times 10/\num{9e16} =
\qty{-1.1e-11}{C/m}$：每米约
$7 \times 10^7$ 个多余电子。电子流在此参考系中更快，是较密的那一束
——每束流各自按自身的 $\gamma$ 收缩。
\textbf{10.} $E' = 2k|\lambda'|/d = 2 \times \num{9e9} \times
\num{1.1e-11}/0.01 = \qty{20}{V/m}$，指向导线；力
$qE' = \qty{2e-5}{N}$，为吸引力。
\textbf{11.} 与 $qvB$ 相同，仅相差因子 $\gamma_v = 1 +
\num{5.6e-8}$：恰为横向力的变换律
（四维力），$F_{\text{静止}} = \gamma F_{\text{实验室}}$。
\textbf{12.} 导线仍载有（更大的）电流，因此 $\vect
B' \neq 0$——但此处我们的电荷静止，而磁场
只作用于运动电荷。
\textbf{13.} $|\lambda'|/\lambda = \gamma_v vu/c^2 = \num{8.2e-16}$.
\textbf{14.} 它乘以 $\lambda/e \approx \num{8.5e22}$
每米元电荷数：$\num{8.2e-16} \times \num{8.5e22}
\approx \num{7e7}$ 个电子（每米）——宏观量级。物质带电如此之多，
其电中性竟如刀刃上的平衡；相对论
则将这刀刃推向一侧。
\textbf{15.} $F/L = \mu_0I_1I_2/2\pi d = \num{2e-7} \times 100/0.01 =
\qty{2e-3}{N/m}$——曾用来\emph{定义}安培的公式。
\textbf{16.} $B = \mu_0nI = 4\pi\times10^{-7} \times \num{e4} \times
10 = \qty{0.13}{T}$；磁压强 $B^2/2\mu_0 \approx
\qty{6300}{Pa} \approx \qty{0.6}{N/cm^2}$：整车的重量悬于
整合后的相对论之上。
\textbf{17.} 令 $c \to \infty$，$\lambda'$ 消失，连同其力，
场的磁分量、电动机、发电机与废料场起重机：
磁现象没有非相对论的存在形式。
\textbf{18.} 静止时，电荷看到的是实验室中电中性的导线：两束
流相互抵消，抵消的精度即物质呈电中性的精度——
经实验已知精度极高，故无力作用。
\textbf{19.} “给定源”的图景只是某一参考系下
单一张量的一个切片：$\vect B$ 是
电磁场中由给定观察者的运动赋予
电流而非电荷的那一部分。
\textbf{20.} 电子自旋：每个电子都是一个基本磁铁
（一种内禀的、量子的“电流”），而铁是这些磁铁
取向一致的物质——在\cref{ch:b3:spin-two-level}与
\cref{ch:b3:electromagnetism-in-matter}中讨论此事。
\textbf{21.} 在束流的静止参考系中，斥力为纯库仑力，
且\emph{最大}；在实验室中，平行电流的磁
吸引几乎将其抵消。两者并无矛盾：实验室中的横向力
即静止参考系中的横向力除以 $\gamma$，而较慢的动力学
（时间延缓）补全了这一解释。
\textbf{22.} 实验室中的净力 $= qE_{\text{实验室}} - qvB_{\text{实验室}} =
(1 - \beta^2)qE_{\text{实验室}} = F_{\text{库仑}}/\gamma^2$，依变换关系；
在静止系中此因子一目了然：纯
静电学，其发散须借延缓的时间观察。
\textbf{23.} 坡印廷矢量：能量以接近 $c$ 的速度穿过导体周围的
\emph{场}，电子
仅起调度作用；电线中的铜是引导，而非管道。
\textbf{24.} $vu/c^2 \approx \num{8e-21}$——然而奥斯特于 1820 年
看到他的指南针在导线旁摆动：这一
$10^{22}$ 个电荷的放大因子，早在一个世纪前就已发挥作用，
而当时无人能为其命名。
\textbf{25.} 一根电中性的 \qty{10}{A} 导线，在以 \qty{e5}{m/s} 运动的参考系中观察，
带有 $\qty{-1.1e-11}{C/m}$：长度收缩的 $10^{-16}$
核算差异，乘以每米 $10^{23}$ 个电荷，
\emph{正是}磁力——磁现象的全部，即相对论
以蜗牛般的速度加以审计。
\end{solution}
